\documentclass[brasileiro,11pt]{letter}
\usepackage[T1]{fontenc}
\usepackage[latin1]{inputenc}
\usepackage{babel}
\usepackage[top=.5in,bottom=.5in,left=1in,right=1in,noheadfoot]{geometry}
\usepackage{graphicx}

\address{\sffamily
  \begin{tabular}{rr}
    Prof.~Dr.~Ant�nio Augusto Fr�hlich & \raisebox{-0.9cm}[0pt][0pt]{\resizebox*{!}{1.2cm}{\includegraphics{/home/guto/etc/logos/lisha}}}\\
    Universidade Federal de Santa Catarina &\\
    Laborat�rio de Integra��o de Software e Hardware  &\\
    \multicolumn{2}{r}{CP 476, 88040-900 Florian�polis, SC}\\
    \multicolumn{2}{r}{Phone: +48 3721-9516}\\
    \multicolumn{2}{r}{E-Mail: guto@lisha.ufsc.br}\\
    \multicolumn{2}{r}{WWW: http://lisha.ufsc.br/$\sim$guto}\\
    \\
    \multicolumn{2}{r}{\today}
  \end{tabular}
}

\date{}

\signature{
  \begin{center}
   \raisebox{-.9cm}[0pt][0pt]{\includegraphics[scale=1]{/home/guto/personal/signature}}\\
    \hrulefill\\
    Ant�nio Augusto Fr�hlich
  \end{center}
}


\begin{document}

\begin{letter}{Revista de Ci�ncia \& Tecnologia\\
    Presidente da Comiss�o Cient�fico-Editorial\
    Prof.~Dr.~Valmir Eduardo Alcarde}
\opening{Prezado Prof. Alcarde,}
\setlength{\parskip}{1.5ex}

Encaminhamos em anexo a vers�o revisada do artigo ``Interactive Digital
TV for Developing Countries: a software/hardware perspective'', de
autoria de Ant�nio Augusto Fr�hlich e Valter Roesler. 

Antes de comentarmos sobre os pareceres dos revisores, gostar�amos de
agradecer a dedica��o e o empenho dedicados ao artigo. As recomenda��es
foram todas atendidas, conforme pode ser visto na vers�o com altera��es
destacadas. Em particular, o coment�rio do primeiro revisor sobre a
forma emp�rica como as quest�es de desempenho da linguagem JAVA s�o
tratados no artigo nos levou a incluir uma refer�ncia a um artigo que
reporta an�lises detalhadas do desempenho desta em sistemas embarcados
de tempo real. J� a sugest�o do segundo revisor sobre o middleware GINGA
foi abordada com a adi��o de informa��es sobre seus componentes e
contexto de utiliza��o.

Atenciosamente,

\closing{~}

\end{letter}

\end{document}





















































